% 第二章示例：说明绪论各部分应回答的问题。
% 正式写作时应结合实际课题替换内容，文件名不影响章节编号。
\chapter{绪论}

\section{研究背景和意义}

绪论用于说明研究从何而来、为什么值得开展，以及论文准备解决什么
问题。研究背景应从所属技术领域、行业需求或理论问题逐步收束到具体
课题，避免堆砌与研究目标关系不大的宏观材料。研究意义可以分别说明
理论价值和实践价值，但应建立在真实问题和已有研究基础上，不宜使用
空泛口号或夸大性表述。

对于毕业设计类课题，还应交代应用场景、设计需求、主要约束和预期
成果；对于实证研究类课题，则应说明研究对象、关键现象和需要验证的
问题。背景、意义和后文研究目标之间应形成清楚的逻辑关系。

\section{国内外研究现状}

研究现状不是参考文献的逐篇摘要，而是围绕研究问题对已有成果进行
分类、比较和评价。作者可按照理论观点、研究方法、技术路线、应用
场景或时间发展过程组织文献，说明不同研究之间的联系、差异和不足。
引用他人观点、数据或方法时，应在正文相应位置标注文献来源
\cite{liu1988vortex,wu2010film,sui2015edm}。

本节结尾通常需要在文献分析基础上归纳尚未解决的问题，说明本文的
切入点。评价已有研究时应保持客观，既不能只罗列优点，也不能在缺少
证据的情况下笼统声称“国内研究较少”或“现有方法均存在缺陷”。

\section{研究目标与主要内容}

研究目标应与题目、研究问题和最终结论保持一致，并尽量使用可以检验
或评价的表述。主要内容用于说明为实现目标将完成哪些相互关联的工作，
而不是简单重复章节名称。一般可以从以下方面组织：

\begin{enumerate}
  \item 明确研究对象、问题边界、需求条件和评价指标；
  \item 说明采用的理论、模型、技术方案、实验方法或调查方法；
  \item 完成数据获取、系统实现、实验验证或案例分析；
  \item 对结果进行比较和解释，归纳结论、创新点与适用范围。
\end{enumerate}

若论文包含创新点，应准确说明创新发生在研究对象、方法、数据、
技术实现还是应用场景，避免把常规工作包装为创新。

\section{论文组织结构}

论文组织结构应简要说明各章承担的任务以及章与章之间的逻辑关系。
正式论文通常从绪论开始，随后依次展开理论基础或相关技术、研究设计
与实施、结果分析与讨论，最后给出总结与展望。不同专业可以根据课题
类型调整章节名称和顺序。

当前示例的第一章是模板使用教程，正式提交前可以删除。从本章开始，
后续章节依次演示研究设计、结果表达、讨论和总结等常见内容。组织
结构说明应随最终目录同步修改，不应保留与实际章节不一致的占位文字。
